%%-----------------------------------%%
%%   MODELAGEM ENTIDADES SQL - 3FN   %%
%%-----------------------------------%%
\section{Modelagem Entidades SQL - 3FN}\label{sec:modelagem}

\begin{explicacao}[Como ler as caixas de entidade]
Cada caixa a seguir representa uma \textbf{tabela do PostgreSQL} na Terceira Forma Normal (3FN). Todos os atributos são \textbf{NOT NULL} por padrão, salvo quando explicitamente indicado \chave{NULL}. O nome do campo aparece à esquerda e o \textbf{tipo de dado} sugerido para o PostgreSQL aparece à direita, em cinza. Quando aplicável, um selo adicional indica \chave{PK} (chave primária), \chave{FK} (chave estrangeira) ou uma restrição como \chave{NULL} (aceita valor nulo) ou \chave{UNIQUE}.
\end{explicacao}

%%-----------------------------------%%
%%        ENTIDADES - USUÁRIO        %%
%%-----------------------------------%%
\subsection{Usuario}
\begin{entidade}{Usuario}
\campo{id}{SERIAL}{PK}
\campo{nome}{VARCHAR(150)}{NOT NULL}
\campo{email}{VARCHAR(150)}{UNIQUE, NOT NULL}
\campo{profile\_url\_photo}{TEXT}{NULL}
\end{entidade}

\subsection{Chefe Geral}
\begin{entidade}{Chefe\_Geral}
\campo{id\_usuario}{INTEGER}{PK, FK}
\end{entidade}

\subsection{Gestor de Almoxarifado}
\begin{entidade}{Gestor\_Almoxarifado}
\campo{id\_usuario}{INTEGER}{PK, FK}
\end{entidade}

\subsection{Tabela Gestor de Almoxarifado \texorpdfstring{\texorpdfstring{$\times$}{x}}{x} Almoxarifado (N para N)}
\begin{entidade}{Gestor\_Almoxarifado\_x\_Almoxarifado}
\campo{id\_gestor\_almoxarifado}{INTEGER}{PK, FK}
\campo{id\_almoxarifado}{INTEGER}{PK, FK}
\end{entidade}
\begin{regra}
A chave primária composta \textit{(id\_gestor\_almoxarifado, id\_almoxarifado)} impede vincular o mesmo gestor duas vezes ao mesmo almoxarifado.
\end{regra}

\subsection{Gestor de bens patrimoniais}
\begin{entidade}{Gestor\_Bens\_Patrimoniais}
\campo{id\_usuario}{INTEGER}{PK, FK}
\end{entidade}
%%-----------------------------------%%
%%        BENS PATRIMONIAIS          %%
%%-----------------------------------%%
\subsection{Bem patrimonial}
\begin{entidade}{Bem\_Patrimonial}
\campo{id}{SERIAL}{PK}
\campo{id\_resumo\_bem\_patrimonial}{INTEGER}{FK, NOT NULL}
\campo{numero\_patrimonio}{VARCHAR(30)}{UNIQUE, NOT NULL}
\campo{estado\_conservacao}{ENUM}{BOM, REGULAR, RUIM, NOT NULL}
\campo{id\_local}{INTEGER}{FK, NOT NULL}
\campo{photo\_url}{TEXT}{NOT NULL}
\campo{documento\_dado\_baixa\_pdf\_url}{TEXT}{NULL}
\campo{nome\_responsavel\_sei}{VARCHAR(150)}{NOT NULL}
\campo{status}{ENUM}{Ativo, Inservivel, Ja\_dado\_baixa, NOT NULL}
\campo{descricao\_complementar}{VARCHAR(200)}{NULL}
\campo{versao}{INTEGER}{DEFAULT 1, NOT NULL}
\end{entidade}

\begin{explicacao}
\textit{id\_resumo\_bem\_patrimonial}, \textit{numero\_patrimonio}, \textit{estado\_conservacao}, \textit{id\_local}, \textit{photo\_url}, \textit{nome\_responsavel\_sei} e \textit{status} são estritamente \textbf{NOT NULL}. Apenas \textit{documento\_dado\_baixa\_pdf\_url} admite NULL (preenchido exclusivamente quando o bem atinge o status \texttt{Ja\_dado\_baixa}) e \textit{descricao\_complementar} (opcional para notas individuais daquela unidade física).
\end{explicacao}

\subsection{Resumo bem patrimonial}
\begin{entidade}{Resumo\_Bem\_Patrimonial}
\campo{id}{SERIAL}{PK}
\campo{nome}{VARCHAR(150)}{NOT NULL}
\campo{descricao}{TEXT}{NOT NULL}
\end{entidade}

\begin{explicacao}[Regra de Identidade de Bens Patrimoniais]
\textit{Resumo\_Bem\_Patrimonial} representa o modelo catalográfico do equipamento (ex.: ``Microscópio Óptico Binocular Coleman''). Vários bens físicos compartilham o mesmo resumo. A alteração do nome do resumo é restrita ao Gestor de Bens e reflete em todos os itens do modelo. Caso um professor solicite alteração de nome para um bem individual, a operação representa reclassificação para outro resumo existente ou criação de um novo modelo, preservando a integridade dos demais equipamentos. Expressamente, resumos de catálogo não possuem vínculo com arquivos fotográficos, pertencendo a imagem unicamente à unidade física cadastrada em \textit{Bem\_Patrimonial}.
\end{explicacao}
\newpage
\subsection{Historico bem patrimonial}
\begin{entidade}{Historico\_Bem\_Patrimonial}
\campo{id}{SERIAL}{PK}
\campo{id\_bem\_patrimonial}{INTEGER}{FK, NOT NULL}
\campo{timestamp}{TIMESTAMP}{NOT NULL}
\campo{tipo}{ENUM}{cadastro, edicao, baixa, NOT NULL}
\campo{id\_usuario}{INTEGER}{FK, NOT NULL}
\end{entidade}

\subsection{Alteracao bem patrimonial}
\begin{entidade}{Alteracao\_Bem\_Patrimonial}
\campo{id}{SERIAL}{PK}
\campo{id\_historico}{INTEGER}{FK, NOT NULL}
\campo{campo}{VARCHAR(60)}{NOT NULL}
\campo{valor\_anterior}{TEXT}{NOT NULL}
\campo{valor\_novo}{TEXT}{NOT NULL}
\end{entidade}
\newpage
%%-----------------------------------%%
%%        REQUISIÇÕES DE BEM         %%
%%-----------------------------------%%
\subsection{Requisicao Professor edicao Bem patrimonial}
\begin{entidade}{Requisicao\_Edicao\_Bem\_Patrimonial}
\campo{id}{SERIAL}{PK}
\campo{id\_bem\_patrimonial}{INTEGER}{FK, NOT NULL}
\campo{status}{ENUM}{pendente, aprovada, rejeitada, NOT NULL}
\campo{feita\_em}{TIMESTAMP}{NOT NULL}
\campo{respondida\_em}{TIMESTAMP}{NULL}
\campo{id\_usuario\_solicitante}{INTEGER}{FK, NOT NULL}
\campo{id\_usuario\_respondente}{INTEGER}{FK, NULL}
\campo{novo\_nome}{VARCHAR(150)}{NULL}
\campo{novo\_status}{ENUM}{Ativo, Inservivel, NULL}
\campo{novo\_estado\_conservacao}{ENUM}{BOM, REGULAR, RUIM, NULL}
\campo{novo\_id\_local}{INTEGER}{FK, NULL}
\campo{nova\_photo\_url}{TEXT}{NULL}
\campo{motivo}{TEXT}{NOT NULL}
\campo{justificativa\_resposta}{TEXT}{NULL}
\campo{novo\_id\_resumo\_bem\_patrimonial}{INTEGER}{FK, NULL}
\campo{versao\_bem\_origem}{INTEGER}{NOT NULL}
\end{entidade}

\begin{regra}
RF10: \textit{UNIQUE(id\_bem\_patrimonial) WHERE status = 'pendente'} -- impede mais de uma requisição de edição pendente simultânea para o mesmo bem patrimonial. Um índice único parcial resolve isso no PostgreSQL; no Firestore, a checagem precisa ser feita na própria \textit{Cloud Function} de criação, dentro de uma transação (ver seção \textit{Tecnologia e Relatórios}).
\end{regra}
\newpage
\subsection{Requisicao Professor adicao Bem patrimonial}
\begin{entidade}{Requisicao\_Adicao\_Bem\_Patrimonial}
\campo{id}{SERIAL}{PK}
\campo{numero\_patrimonio\_proposto}{VARCHAR(30)}{NOT NULL}
\campo{status}{ENUM}{pendente, aprovada, rejeitada, NOT NULL}
\campo{feita\_em}{TIMESTAMP}{NOT NULL}
\campo{id\_bem\_patrimonial\_se\_aprovado}{INTEGER}{FK, NULL}
\campo{respondida\_em}{TIMESTAMP}{NULL}
\campo{id\_usuario\_solicitante}{INTEGER}{FK, NOT NULL}
\campo{id\_usuario\_respondente}{INTEGER}{FK, NULL}
\campo{estado\_conservacao\_proposto}{ENUM}{BOM, REGULAR, RUIM, NOT NULL}
\campo{photo\_url\_proposta}{TEXT}{NOT NULL}
\campo{nome\_responsavel\_proposto}{VARCHAR(150)}{NOT NULL}
\campo{id\_local}{INTEGER}{FK, NOT NULL}
\campo{id\_resumo\_bem\_patrimonial}{INTEGER}{FK, NULL}
\campo{nome\_resumo\_proposto}{VARCHAR(150)}{NULL}
\campo{descricao\_resumo\_proposta}{TEXT}{NULL}
\campo{motivo}{TEXT}{NOT NULL}
\campo{justificativa\_resposta}{TEXT}{NULL}
\end{entidade}

\begin{explicacao}
Se o local desejado não existir no banco, o professor deve primeiro criar o \textit{Local} no sistema, para então vincular a respectiva FK (\textit{id\_local}) nesta requisição.
\end{explicacao}
\begin{regra}
RF10b: \textit{UNIQUE(numero\_patrimonio\_proposto) WHERE status = 'pendente'} --
impede mais de uma requisição de adição pendente simultânea para o mesmo
número de patrimônio. No PostgreSQL, isso é implementado como índice
único parcial; no Firestore, a verificação é reimplementada na Cloud
Function de criação da requisição, dentro de uma transação.
\end{regra}
\subsection{Lock de Requisição de Patrimônio}
\begin{entidade}{Locks\_Requisicao\_Patrimonio}
\campo{id (docId determinístico)}{VARCHAR(100)}{PK}
\campo{criado\_em}{TIMESTAMP}{NOT NULL}
\end{entidade}

\begin{regra}[Lock Determinístico para Unicidade de Requisição Pendente]
O Firestore não adquire trava pessimista sobre \textbf{ausência} de documento em uma query transacional. Por isso, a garantia de que não existam duas requisições pendentes para o mesmo bem ou plaqueta \textbf{não pode} ser implementada pela consulta \texttt{WHERE status = 'pendente'} dentro de uma transação -- dois clientes simultâneos lerão conjunto vazio, ambos passarão na validação e ambos gravarão, violando RF10/RF10b.

A solução é o padrão \textit{deterministic lock}: a Cloud Function tenta criar um documento cujo \textbf{docId é derivado deterministicamente} do recurso disputado. O Firestore garante que apenas uma gravação com o mesmo docId vence quando executadas concorrentemente dentro de uma transação.

\begin{itemize}
  \item Para requisição de \textbf{edição}: docId = \texttt{bem\_edicao\_\{id\_bem\_patrimonial\}}
  \item Para requisição de \textbf{adição}: docId = \texttt{bem\_adicao\_\{numero\_patrimonio\_proposto\}}
\end{itemize}

O lock é removido na transação de aprovação/rejeição ordinária. O lock \texttt{Locks\_Requisicao\_Patrimonio/bem\_edicao\_\{id\}} DEVE ser removido pela transação também em caso de rejeição por divergência de versão do bem patrimonial, solucionando deadlocks permanentes. Lock e requisição são criados atomicamente; não há commit parcial dessa transação. Limpeza verifica o vínculo e jamais remove um lock apenas porque a requisição pendente tem mais de 24 horas.
\end{regra}

\subsection{Impressão de Etiquetas de Frasco}

\begin{entidade}{Impressao\_Etiqueta\_Frasco}
\campo{id}{SERIAL}{PK}
\campo{gerado\_em}{TIMESTAMP}{NOT NULL}
\campo{gerado\_por}{INTEGER}{FK, NOT NULL}
\campo{codigo\_inicial}{INTEGER}{NOT NULL}
\campo{codigo\_final}{INTEGER}{NOT NULL}
\end{entidade}

\begin{explicacao}
Esta entidade audita de forma isolada as sessões de geração em lote de etiquetas virgens contínuas. \textit{codigo\_inicial} e \textit{codigo\_final} referem-se aos números sequenciais gerados (ex.: \texttt{LCQUI-1} a \texttt{LCQUI-50}). Não há vínculo de chave estrangeira com a tabela \textit{Frasco\_Reagente}, dado que a impressão em branco é puramente utilitária e não cria reservas de ID. A transação e consolidação final do código só ocorrem no banco no momento do cadastro físico.
\end{explicacao}

%%-----------------------------------%%
%%        ALMOXARIFADO / REAGENTES   %%
%%-----------------------------------%%
\subsection{Almoxarifado}
\begin{entidade}{Almoxarifado}
\campo{id}{SERIAL}{PK}
\campo{id\_local}{INTEGER}{FK, NOT NULL}
\campo{nome}{VARCHAR(100)}{NOT NULL}
\campo{descricao}{VARCHAR(500)}{NOT NULL}
\campo{ativo}{BOOLEAN}{DEFAULT TRUE, NOT NULL}
\end{entidade}

\begin{explicacao}
O campo \textit{ativo} indica se o almoxarifado está operacional. Quando \texttt{ativo = false}, o almoxarifado está desativado e as regras de bloqueio da Seção~\ref{sec:regras-almoxarifado-inativo} se aplicam. A desativação é reversível e não exclui registros históricos.
\end{explicacao}

\subsection{Substância Química}
\begin{entidade}{Substancia\_Quimica}
\campo{id}{SERIAL}{PK}
\campo{nome}{VARCHAR(150)}{NOT NULL}
\campo{cas\_number}{VARCHAR(15)}{UNIQUE, NULL}
\campo{formula\_quimica}{VARCHAR(50)}{NULL}
\campo{ativo}{BOOLEAN}{DEFAULT TRUE, NOT NULL}
\end{entidade}

\begin{explicacao}
A entidade \textit{Substancia\_Quimica} representa uma substância quimicamente identificável, independentemente de fabricante, concentração, pureza, lote ou frasco.

O campo \textit{cas\_number} identifica a substância e permanece \textbf{UNIQUE} (permitindo múltiplos NULL, já que algumas substâncias podem não ter CAS cadastrado). Uma mesma substância pode participar da composição de diversas especificações de reagentes diferentes.
\end{explicacao}
\newpage
\subsection{Resumo do Reagente / Solvente}
\begin{entidade}{Resumo\_Reagente}
\campo{id}{SERIAL}{PK}
\campo{nome}{VARCHAR(150)}{NOT NULL}
\campo{tipo\_substancia}{ENUM}{PURA, MISTURA, NOT NULL}
\campo{natureza\_quimica}{ENUM}{ORGANICO, INORGANICO, ELEMENTO, HIBRIDO, NOT NULL}
\campo{requer\_pesagem\_frequente}{BOOLEAN}{DEFAULT FALSE, NOT NULL}
\campo{frequencia\_pesagem\_dias}{INTEGER}{NULL}
\campo{estado\_fisico}{ENUM}{SOLIDO, LIQUIDO, NOT NULL}
\campo{eh\_higroscopico}{BOOLEAN}{DEFAULT FALSE, NOT NULL}
\end{entidade}
\newpage
\begin{explicacao}[Regra Fechada de Identidade do Resumo de Reagente]
O \textit{Resumo\_Reagente} representa o item químico agrupador apresentado no catálogo do sistema, permitindo agrupar diferentes especificações comerciais ou formulações do mesmo reagente.

Dois reagentes pertencem ao \textbf{mesmo} \textit{Resumo\_Reagente} quando compartilham a mesma identidade química nominal de catálogo (ex.: ``Ácido Clorídrico'', ``Etanol'', ``Hidróxido de Sódio'').
\begin{itemize}
    \item Variações de pureza (P.A., Técnico, Grau HPLC) $\implies$ \textbf{mesmo resumo}, especificações distintas.
    \item Variações de concentração (HCl 37\%, HCl 10\%, Etanol 99.5\%, Etanol 70\%) $\implies$ \textbf{mesmo resumo}, especificações distintas.
    \item Fabricante, código de catálogo e densidade $\implies$ atributos de especificação, \textbf{mesmo resumo}.
\end{itemize}
Devem ser cadastrados como \textbf{resumos diferentes}:
\begin{itemize}
    \item Espécies químicas distintas com fórmulas ou números CAS bases diferentes (ex.: Sulfato de Cobre II vs Sulfato de Ferro II);
    \item Soluções formuladas consagradas com finalidade analítica própria (ex.: ``Reagente de Benedict'', ``Solução de Fehling A'', ``Solução de Lugol'') $\implies$ cada uma constitui um \textit{Resumo\_Reagente} próprio do tipo \texttt{MISTURA}.
\end{itemize}

Os campos \textit{medida\_total\_usado} e \textit{unidade\_de\_medida} foram removidos desta entidade: ambos são métricas agregadas derivadas de \textit{Frasco\_Reagente} e \textit{Emprestimo\_Reagente}, e devem existir apenas na view materializada descrita na seção \textit{Métricas Agregadas}, nunca como coluna editável de uma tabela cadastral.

A unidade operacional é derivada de Resumo\_Reagente.estado\_fisico: SOLIDO usa g e LIQUIDO usa mL. Não persistir cópia redundante da unidade no modelo 3FN.

DP-A01: estado\_fisico e eh\_higroscopico pertencem ao Resumo\_Reagente. Estados físicos distintos nas condições de catálogo (25 graus Celsius, 1 atm) exigem resumos distintos; solução aquosa e sólido puro não são o mesmo resumo. Densidade e concentração permanecem na especificação comercial.

\textit{natureza\_quimica} classifica o item do catálogo como um todo (não uma \textit{Substancia\_Quimica} individual): \texttt{ELEMENTO} para substâncias simples (não se classificam como orgânicas nem inorgânicas), \texttt{ORGANICO}/\texttt{INORGANICO} para os demais casos, e \texttt{HIBRIDO} para fluidos complexos que não se decompõem em uma lista simples de substâncias em \textit{Composicao\_Reagente} (ex.: Soro Fetal Bovino). Como \texttt{HIBRIDO} não é derivável a partir das substâncias componentes, o campo é curado no nível do resumo, e não em \textit{Substancia\_Quimica}, evitando duas fontes de verdade diferentes para itens PUROS e para MISTURA/HÍBRIDO. O rótulo exibido na interface (ex.: ``Híbrido'', ``Complexo'' ou ``Biológico'') é independente do valor salvo no banco e pode ser alterado livremente.

O campo \textit{letra\_inicial}, presente em versões anteriores deste documento, foi removido do modelo 3FN: por ser derivável do próprio \textit{nome}, sua persistência aqui violaria a forma normal. Ele reaparece como uma denormalização proposital exclusiva da implementação em Firestore -- ver seção \textit{Notas de Mapeamento para Firestore}.
\end{explicacao}
\newpage
\subsection{Composição do Reagente}
\begin{entidade}{Composicao\_Reagente}
\campo{id}{SERIAL}{PK}
\campo{id\_especificacao\_reagente}{INTEGER}{FK, NOT NULL}
\campo{id\_substancia\_quimica}{INTEGER}{FK, NOT NULL}
\campo{valor\_min}{NUMERIC(10,5)}{NULL}
\campo{valor\_max}{NUMERIC(10,5)}{NULL}
\campo{tipo\_concentracao}{ENUM}{M\_M,V\_V,M\_V,MOL\_L,MOL\_KG,PPM,PPB}
\campo{notacao\_original\_fabricante}{VARCHAR(50)}{NULL}
\end{entidade}

\begin{explicacao}
A entidade \textit{Composicao\_Reagente} representa a relação entre uma especificação de reagente cadastrada e as substâncias químicas que a compõem.

A relação é do tipo N:N:

\begin{center}
\textit{Especificacao\_Reagente} $\leftrightarrow$ \textit{Substancia\_Quimica}
\end{center}

A composição foi movida do nível de \textit{Resumo\_Reagente} para o nível de \textit{Especificacao\_Reagente}, pois a concentração (37\%, 10\%, 1 mol/L etc.) é uma característica que varia entre especificações de um mesmo resumo, não uma propriedade do resumo em si. Vincular a composição ao resumo impediria registrar corretamente, por exemplo, HCl 37\% e HCl 10\% sob o mesmo \textit{Resumo\_Reagente} com valores de composição diferentes.

Uma substância pura normalmente terá apenas uma linha associada, enquanto uma mistura poderá possuir diversas substâncias.

Q03/DP-A03: concentração pontual usa valor\_min = valor\_max ou valor\_max NULL; faixa usa limites numéricos ordenados. CHECK (valor\_max IS NULL OR valor\_min <= valor\_max). O servidor rejeita faixa com máximo informado e mínimo ausente; preserva a notação original do fabricante separadamente, sem usá-la como autoridade numérica.
\end{explicacao}

\begin{regra}
\textit{UNIQUE(id\_especificacao\_reagente, id\_substancia\_quimica)} -- impede duas linhas de composição para a mesma substância dentro da mesma especificação.
\end{regra}

\subsection{Especificação do Reagente / Solvente}
\begin{entidade}{Especificacao\_Reagente}
\campo{id}{SERIAL}{PK}
\campo{id\_resumo\_reagente}{INTEGER}{FK, NOT NULL}
\campo{descricao}{VARCHAR(150)}{NOT NULL}
\campo{fabricante}{VARCHAR(150)}{NULL}
\campo{codigo\_produto\_fabricante}{VARCHAR(100)}{NULL}
\campo{grau\_pureza}{VARCHAR(30)}{NULL}
\campo{densidade}{NUMERIC(8,4)}{NULL}
\campo{classe\_inflamabilidade}{ENUM}{NAO\_INFLAMAVEL,CLASSE\_1,CLASSE\_2,CLASSE\_3, NOT NULL}
\campo{eh\_controlado\_pf}{BOOLEAN}{DEFAULT FALSE, NOT NULL}
\campo{eh\_controlado\_eb}{BOOLEAN}{DEFAULT FALSE, NOT NULL}
\campo{link\_fds\_fispq}{VARCHAR(255)}{NULL}
\end{entidade}

\begin{explicacao}
A especificação referencia o resumo, fonte do estado físico; a unidade é derivada desse estado. Frasco e empréstimo chegam ao resumo pela cadeia de FKs. Snapshots de consulta e de cálculo são documentados separadamente no Firestore.

\textit{GASOSO} foi removido de \textit{estado\_fisico} nesta versão: nenhum reagente gasoso está listado para o almoxarifado atual, e reagentes gasosos precisam de campos e método de medição próprios (pressão, não peso em balança) em vez de reaproveitar os campos já pensados para sólido/líquido. O estudo detalhado para reintrodução está na seção \textit{Implementações em Estudo para Versões Futuras}.
\end{explicacao}

\subsection{Estoque Mínimo por Almoxarifado}
\begin{entidade}{Estoque\_Minimo\_Almoxarifado}
\campo{id\_especificacao\_reagente}{INTEGER}{FK, PK, NOT NULL}
\campo{id\_almoxarifado}{INTEGER}{FK, PK, NOT NULL}
\campo{qtd\_limiar\_escassez}{INTEGER}{NOT NULL}
\campo{ativo}{BOOLEAN}{DEFAULT TRUE, NOT NULL}
\campo{notificacao\_ativa}{BOOLEAN}{DEFAULT TRUE, NOT NULL}
\end{entidade}

\begin{explicacao}
A entidade \textit{Estoque\_Minimo\_Almoxarifado} estabelece o limiar de escassez (em quantidade de frascos) de uma \textit{Especificacao\_Reagente} para um determinado \textit{Almoxarifado}. O par Especificação × Almoxarifado é a identidade unívoca da configuração. No Firestore, esta entidade é implementada como a subcoleção \texttt{Almoxarifado/\{id\}/Estoques\_Configurados/\{idResumo\_idEspec\}}.
O campo \textit{ativo} indica que o par continua configurado ou intencionalmente estocado. 
O campo \textit{notificacao\_ativa} indica que a emissão automática do alerta está habilitada (útil para silenciar alertas sem excluir a intenção de estocar a especificação).
\end{explicacao}

\subsection{Lote}
\begin{entidade}{Lote}
\campo{id}{SERIAL}{PK}
\campo{id\_especificacao\_reagente}{INTEGER}{FK, NOT NULL}
\campo{data\_aquisicao}{TIMESTAMP}{NOT NULL}
\campo{qtd\_frascos\_comprados}{INTEGER}{DEFAULT 0, NOT NULL}
\campo{nome\_fornecedor}{VARCHAR(80)}{NOT NULL}
\campo{numero\_lote}{VARCHAR(100)}{NOT NULL}
\campo{nota\_fiscal}{VARCHAR(100)}{NOT NULL}
\campo{data\_fabricacao}{DATE}{NOT NULL}
\campo{data\_validade}{DATE}{NOT NULL}
\end{entidade}
\begin{regra}
\textit{UNIQUE(id\_especificacao\_reagente, nome\_fornecedor, numero\_lote)}.
\end{regra}
\begin{explicacao}
\textit{qtd\_frascos\_comprados} é dado de entrada (o que consta na nota fiscal) e permanece nesta tabela. Já \textit{qtd\_frascos\_cadastrados} -- quantos frascos deste lote já foram individualmente registrados em \textit{Frasco\_Reagente} -- é uma contagem agregada (\texttt{COUNT(*) FROM Frasco\_Reagente WHERE id\_lote = X}) e, por isso, \textbf{não} é coluna desta tabela cadastral: ela vive exclusivamente em \textit{Lote\_Materializado} (ver seção \textit{Materialized (Views)}), mantida por mecanismo transacional a cada frasco cadastrado ou removido.
\end{explicacao}
\subsection{Frasco de Reagente / Solvente}
\begin{entidade}{Frasco\_Reagente}
\campo{id}{SERIAL}{PK}
\campo{id\_almoxarifado}{INTEGER}{FK, NOT NULL}
\campo{id\_lote}{INTEGER}{FK, NULL}
\campo{id\_especificacao\_reagente}{INTEGER}{FK, NULL}
\campo{detalhe\_local\_armazenamento}{TEXT}{NULL}
\campo{codigo\_frasco}{TEXT}{UNIQUE, NOT NULL}
\campo{data\_ultima\_pesagem}{TIMESTAMP}{NULL}
\campo{conteudo\_nominal}{NUMERIC(10,3)}{NULL}
\campo{peso\_no\_cadastrado}{NUMERIC(10,3)}{NOT NULL}
\campo{peso\_atual}{NUMERIC(10,3)}{NOT NULL}
\campo{peso\_frasco\_vazio}{NUMERIC(10,3)}{NULL}
\campo{medida\_usada}{NUMERIC(10,3)}{DEFAULT 0, NOT NULL}
\campo{data\_abertura}{DATE}{NULL}
\campo{validade\_fechado}{DATE}{NULL}
\campo{validade\_efetiva}{DATE}{NULL}
\campo{validade\_desconhecida}{BOOLEAN}{DEFAULT FALSE, NOT NULL}
\campo{validade\_apos\_aberto\_dias}{INTEGER}{NULL}
\campo{estado\_fisico\_frasco}{ENUM}{FECHADO,ABERTO,VAZIO,QUEBRADO,DESCARTADO,EXTRAVIADO, NOT NULL}
\campo{disponibilidade}{ENUM}{DISPONIVEL, EMPRESTADO, INDISPONIVEL, NOT NULL}
\campo{vencido}{BOOLEAN}{DEFAULT FALSE, NOT NULL}
\campo{em\_quarentena}{BOOLEAN}{DEFAULT FALSE, NOT NULL}
\campo{uso\_vencido\_autorizado}{BOOLEAN}{DEFAULT FALSE, NOT NULL}
\campo{detalhe\_status}{TEXT}{NULL}
\campo{cadastrado\_em}{TIMESTAMP}{NOT NULL}
\campo{cadastrado\_por}{INTEGER}{FK, NOT NULL}
\campo{abertura\_historica\_desconhecida}{BOOLEAN}{DEFAULT FALSE, NOT NULL}
\end{entidade}

\begin{regra}[Constraint de Identidade Química do Frasco]
\textit{CHECK (id\_lote IS NOT NULL OR id\_especificacao\_reagente IS NOT NULL)} -- pelo menos um dos dois deve existir. Quando \textit{id\_lote} é informado, \textit{id\_especificacao\_reagente} fica NULL (a especificação é derivável via \textit{Lote.id\_especificacao\_reagente}). Quando \textit{id\_lote} é NULL, \textit{id\_especificacao\_reagente} é obrigatório para manter o caminho até a densidade, unidade de medida e resumo do reagente.
\end{regra}

\begin{explicacao}
\textit{conteudo\_nominal} (antigo \textit{capacidade\_nominal}) representa o conteúdo declarado no cadastro: para líquidos, é o volume nominal em mL; para sólidos, é a massa nominal em g. Se NULL, o valor original do rótulo é desconhecido ou ilegível (não usar zero para representar desconhecido).

\textit{peso\_no\_cadastrado} é o peso total (frasco + conteúdo) registrado no momento do cadastro. Este campo é \textbf{imutável} após o cadastro -- nunca é sobrescrito. \textit{peso\_atual} é o peso total corrente, inicializado com o valor de \textit{peso\_no\_cadastrado} e atualizado a cada devolução de empréstimo ou pesagem de rotina.

Os campos \textit{peso\_no\_cadastrado}, \textit{peso\_atual}, \textit{peso\_frasco\_vazio} e \textit{medida\_usada} são sempre leituras de balança em gramas, independentemente de o conteúdo ser sólido ou líquido, já que a operação real do laboratório consiste em pesar o frasco; o volume, quando aplicável, é sempre derivado por conversão via densidade.

O campo \textit{unidade\_de\_medida}, antes duplicado nesta entidade, foi removido: a unidade do produto (ml ou g) é obtida via \textit{Frasco\_Reagente} $\rightarrow$ \textit{Lote} $\rightarrow$ \textit{Especificacao\_Reagente} (ou diretamente via \textit{id\_especificacao\_reagente} quando o lote é NULL), chegando ao Resumo\_Reagente, cujo estado determina a unidade.

O antigo campo único \textit{status} (com valores combinados como \texttt{VENCIDO\_DISPONIVEL}, \texttt{VENCIDO\_EMPRESTADO}) foi decomposto em dimensões ortogonais -- \textit{estado\_fisico\_frasco}, \textit{disponibilidade}, \textit{vencido} e \textit{em\_quarentena} -- evitando o crescimento combinatorial do enum e simplificando consultas e regras de notificação.

O campo \textit{id\_usuario\_em\_uso} foi removido: quem está com o frasco no momento é obtido consultando \textit{Emprestimo\_Reagente} pelo empréstimo ativo (\textit{status} \texttt{EM\_USO} ou \texttt{ATRASADO}) daquele \textit{id\_frasco\_reagente}.

\textit{validade\_fechado} é mantido de forma independente de \textit{Lote.data\_validade} porque \textit{id\_lote} é opcional: um frasco pode ser cadastrado sem lote associado, e mesmo quando o lote existe, o formulário apenas pré-preenche este campo a partir dele -- o valor permanece editável e não é uma cópia somente-leitura.

O campo \textit{uso\_vencido\_autorizado} fecha a distinção de negócio entre um frasco vencido que permanece apto ao uso didático e um frasco vencido que aguarda descarte.
\end{explicacao}

\subsection{Histórico do Frasco de Reagente / Solvente}
\begin{entidade}{Historico\_Frasco\_Reagente}
\campo{id}{SERIAL}{PK}
\campo{id\_frasco\_reagente}{INTEGER}{FK, NOT NULL}
\campo{id\_almoxarifado}{INTEGER}{FK, NOT NULL}
\campo{id\_gestor}{INTEGER}{FK, NOT NULL}
\subcampo{Para ações automáticas do servidor (jobs de vencimento, atraso), no 3FN usa-se o id\_usuario = 0 (\'SISTEMA LCQUI\'); no Firestore, usa-se a string sentinela \_\_SISTEMA\_\_.}
\campo{tipo}{ENUM}{
  \shortstack[l]{%
    CADASTRO, SAIU, ENTROU, FICOU\_VAZIO, QUEBROU, FOI\_DESCARTADO,\\
    VENCEU, INVENTARIO\_ROTINA, EVAPORACAO, AJUSTE, CONCLUSAO,\\
    ENTROU\_EM\_QUARENTENA, LIBERADO\_QUARENTENA, PENDENTE\_DE\_DESCARTE,\\
    USO\_VENCIDO\_AUTORIZADO, EXTRAVIO, REENCONTRO%
  }
}{NOT NULL}
\campo{campo\_ajustado}{VARCHAR(30)}{NULL}
\campo{id\_emprestimo\_reagente}{INTEGER}{FK, NULL}
\campo{peso\_anterior}{NUMERIC(10,3)}{NULL}
\campo{peso\_novo}{NUMERIC(10,3)}{NULL}
\campo{medida\_ajustada}{NUMERIC(10,3)}{NULL}
\campo{unidade\_medida\_ajustada}{ENUM}{ml,g, NULL}
\campo{motivo}{TEXT}{NULL}
\campo{timestamp}{TIMESTAMP}{NOT NULL}
\end{entidade}

\begin{explicacao}
\textit{peso\_anterior} e \textit{peso\_novo} são sempre leituras de balança em gramas. Já \textit{unidade\_medida\_ajustada} qualifica exclusivamente \textit{medida\_ajustada}, indicando se o ajuste concluído é reportado em massa ou volume. Podem ser NULL caso o evento histórico não seja de natureza quantitativa. O campo \textit{motivo} é genérico e aceita NULL, porém é obrigatório para eventos que exigem justificativa (ex.: EXTRAVIO, REENCONTRO e ajustes manuais). No modelo relacional 3FN, reserva-se a linha canônica \texttt{Usuario.id = 0} com nome \'SISTEMA LCQUI\' para satisfazer a \textit{foreign key} não nula em rotinas automáticas de servidor; no Firestore, o documento recebe a string sentinela \texttt{\_\_SISTEMA\_\_}.
\end{explicacao}

\subsection{Empréstimo de Reagente / Solvente}
\begin{entidade}{Emprestimo\_Reagente}
\campo{id}{SERIAL}{PK}
\campo{id\_frasco\_reagente}{INTEGER}{FK, NOT NULL}
\campo{id\_usuario\_retirou}{INTEGER}{FK, NOT NULL}
\campo{id\_gestor\_retirada}{INTEGER}{FK, NOT NULL}
\campo{status}{ENUM}{EM\_USO,ATRASADO,DEVOLVIDO,DEVOLVIDO\_COM\_ATRASO,ENCERRADO\_EXTRAORDINARIO, NOT NULL}
\campo{data\_retirada}{TIMESTAMP}{NOT NULL}
\campo{data\_devolucao\_prevista}{DATE}{NOT NULL}
\campo{data\_devolucao\_efetuada}{TIMESTAMP}{NULL}
\campo{id\_local\_usado}{INTEGER}{FK, NOT NULL}
\campo{id\_usuario\_devolveu}{INTEGER}{FK, NULL}
\campo{id\_gestor\_devolucao}{INTEGER}{FK, NULL}
\campo{peso\_saida}{NUMERIC(10,3)}{NOT NULL}
\campo{peso\_retorno}{NUMERIC(10,3)}{NULL}
\campo{medida\_utilizada}{NUMERIC(10,3)}{NULL}
\campo{unidade\_medida\_utilizada}{ENUM}{ml,g, NOT NULL}
\campo{peso\_perda\_evaporacao}{NUMERIC(10,3)}{DEFAULT 0, NOT NULL}
\campo{uso\_vencido\_aceito}{BOOLEAN}{DEFAULT FALSE, NOT NULL}
\campo{finalidade\_uso}{ENUM}{NOT NULL}
\subcampo{AULA\_PRATICA, DEMONSTRACAO, PESQUISA\_TCC\_POS, ESTUDO\_DEGRADACAO\_RESIDUOS}
\campo{auto\_atendimento}{BOOLEAN}{DEFAULT FALSE, NOT NULL}
\campo{justificativa\_metodologica}{TEXT}{NULL}
\campo{tcr\_versao}{VARCHAR(100)}{NULL}
\campo{tcr\_aceito\_em}{TIMESTAMP}{NULL}
\campo{tcr\_aceito\_por}{INTEGER}{FK, NULL}
\campo{tcr\_auditoria\_id}{INTEGER}{FK, NULL}
\campo{tipo\_encerramento\_excepcional}{ENUM}{QUEBRA\_ACIDENTAL, EXTRAVIO\_SINISTRO, NULL}
\campo{motivo\_encerramento\_excepcional}{TEXT}{NULL}
\campo{id\_gestor\_encerramento}{INTEGER}{FK, NULL}
\campo{massa\_perda\_estimada\_g}{NUMERIC(10,3)}{NULL}
\end{entidade}

\begin{explicacao}
\textbf{Semântica dos atores da retirada:} \textit{id\_usuario\_retirou} identifica o \textbf{portador físico} -- o Professor ou Bolsista que receberá e usará o reagente. \textit{id\_gestor\_retirada} identifica o \textbf{operador do almoxarifado} -- o Gestor de Almoxarifado (ou Chefe Geral) que registrou a operação no sistema. No autoatendimento, os dois UIDs coincidem e a flag \textit{auto\_atendimento} fica \texttt{true}. Esta separação é essencial para rastreabilidade de auditoria: o relatório de atividade do gestor parte de \textit{id\_gestor\_retirada}, enquanto o painel ``Meus Reagentes'' do professor parte de \textit{id\_usuario\_retirou}.

\textbf{Semântica dos atores da devolução:} \textit{id\_usuario\_devolveu} identifica o portador físico que entrega o frasco ao balcão, enquanto \textit{id\_gestor\_devolucao} identifica o Gestor de Almoxarifado que operou e registrou a devolução na balança.

\textit{peso\_saida} e \textit{peso\_retorno} são sempre leituras de balança em gramas. \textit{unidade\_medida\_utilizada} qualifica exclusivamente \textit{medida\_utilizada} -- o consumo calculado, reportado em massa (g) para sólidos ou em volume (ml) para líquidos, via $V = \Delta m / \rho$.

O campo \textit{data\_devolucao\_prevista} em nível NoSQL físico armazena uma string civil no padrão YYYY-MM-DD (fuso America/Sao\_Paulo), e o vencimento legal encerra-se estritamente às 23:59:59.999 do referido dia civil.
\end{explicacao}



\subsection{Métricas Agregadas do Reagente}
\begin{explicacao}
As seguintes informações devem ser tratadas preferencialmente como \textbf{views materializadas, consultas agregadas ou campos derivados mantidos por mecanismo transacional}:
\begin{itemize}
\item \textit{medida\_total\_usado} e \textit{unidade\_de\_medida} (agregados por \textit{Resumo\_Reagente});
\item \textit{qtd\_frascos\_fechados\_disponiveis};
\item \textit{qtd\_frascos\_abertos\_disponiveis};
\item \textit{qtd\_frascos\_em\_uso};
\item \textit{qtd\_frascos\_vazios};
\item \textit{qtd\_frascos\_descartados};
\item \textit{qtd\_frascos\_vencidos};
\item \textit{qtd\_frascos\_quebrados}.
\end{itemize}
\end{explicacao}

%%-----------------------------------%%
%%        TURMAS / ROTEIROS          %%
%%-----------------------------------%%
\subsection{Turma}
\begin{entidade}{Turma}
\campo{id}{SERIAL}{PK}
\campo{id\_materia}{INTEGER}{FK, NOT NULL}
\campo{id\_professor}{INTEGER}{FK, NOT NULL}
\campo{status}{ENUM}{Ativo, Arquivada, NOT NULL}
\campo{nome\_turma}{VARCHAR(100)}{NOT NULL}
\campo{ano}{INTEGER}{NOT NULL}
\campo{semestre}{SMALLINT}{CHECK (semestre IN (1,2)), NOT NULL}
\campo{capacidade}{INTEGER}{NOT NULL}
\campo{codigo\_turma}{VARCHAR(20)}{UNIQUE, NOT NULL}
\campo{data\_criacao}{TIMESTAMP}{NOT NULL}
\end{entidade}

\begin{explicacao}
Os campos \textit{ano} e \textit{semestre} representam o per\'{i}odo letivo da turma (ex.: 2026/1). Ambos s\~{a}o exibidos na listagem de turmas arquivadas do professor -- eliminando a necessidade de parsing do \textit{nome\_turma} ou infer\^{e}ncia a partir de \textit{data\_criacao}. O modal de cria\c{c}\~{a}o de turma solicita esses dois campos al\'{e}m dos demais.
\end{explicacao}

\begin{regra}[RN-TUR-01: Controle de Capacidade da Turma]
O ingresso de aluno por código de turma exige obrigatoriamente que a quantidade de alunos matriculados ativos seja estritamente menor que a capacidade da turma:
\[
\text{COUNT}(\text{Aluno\_x\_Turma WHERE id\_turma} = X) < \text{Turma.capacidade}
\]
Caso a capacidade seja atingida, o sistema rejeita o ingresso por código. Apenas o professor responsável pode ultrapassar a capacidade mediante envio de convite nominal por e-mail, registrando justificativa em auditoria.
\end{regra}

\subsection{Historico de Alunos na Turma}
\begin{entidade}{Historico\_Alunos\_Turma}
\campo{id}{SERIAL}{PK}
\campo{id\_turma}{INTEGER}{FK, NOT NULL}
\campo{tipo}{ENUM}{inclusao\_aluno, exclusao\_aluno, NOT NULL}
\campo{id\_aluno}{INTEGER}{FK, NOT NULL}
\campo{modo\_ingresso}{ENUM}{CODIGO, CONVITE, NULL}
\campo{justificativa}{TEXT}{NULL}
\campo{timestamp}{TIMESTAMP}{NOT NULL}
\campo{removido\_por}{INTEGER}{FK, NULL}
\subcampo{Obrigatório quando tipo = 'exclusao\_aluno'}
\end{entidade}

\subsection{Historico de Posts na Turma}
\begin{entidade}{Historico\_Posts\_Turma}
\campo{id}{SERIAL}{PK}
\campo{id\_post}{INTEGER}{FK, NOT NULL}
\campo{editado\_por}{INTEGER}{FK, NOT NULL}
\campo{novo\_titulo}{VARCHAR(150)}{NULL}
\campo{novo\_id\_roteiro}{INTEGER}{FK, NULL}
\campo{nova\_descricao}{TEXT}{NULL}
\campo{antigo\_titulo}{VARCHAR(150)}{NULL}
\campo{antigo\_id\_roteiro}{INTEGER}{FK, NULL}
\campo{antiga\_descricao}{TEXT}{NULL}
\campo{timestamp}{TIMESTAMP}{NOT NULL}
\end{entidade}

\begin{explicacao}
O campo \textit{editado\_por} registra o UID do usuário que realizou a edição. Normalmente é o Professor dono da turma, mas o Chefe Geral também pode editar ou remover posts em caráter excepcional de moderação institucional (Q13). O \textit{id\_turma} não é armazenado nesta entidade 3FN porque pode ser obtido por \textit{Historico\_Posts\_Turma.id\_post} $\rightarrow$ \textit{Post.id\_turma}; se a consulta direta por turma for necessária no Firestore, esse vínculo poderá ser denormalizado na coleção de histórico.
\end{explicacao}

\subsection{Historico de Comentario em um Post}
\begin{entidade}{Historico\_Comentario}
\campo{id}{SERIAL}{PK}
\campo{id\_comentario}{INTEGER}{FK, NOT NULL}
\campo{novo\_texto}{TEXT}{NOT NULL}
\campo{texto\_antigo}{TEXT}{NOT NULL}
\campo{editado\_em}{TIMESTAMP}{NOT NULL}
\campo{editado\_por}{INTEGER}{FK, NOT NULL}
\end{entidade}

\subsection{Roteiro de Experimento}
\begin{entidade}{Roteiro\_Experimento}
\campo{id}{SERIAL}{PK}
\campo{id\_professor\_upload}{INTEGER}{FK, NOT NULL}
\campo{file\_url}{TEXT}{NOT NULL}
\campo{nome}{VARCHAR(150)}{NOT NULL}
\campo{descricao}{TEXT}{NOT NULL}
\end{entidade}

\subsection{Roteiro Professor Compartilhado}
\begin{entidade}{Roteiro\_Professor\_Compartilhado}
\campo{id\_roteiro\_experimento}{INTEGER}{PK, FK}
\campo{id\_professor}{INTEGER}{PK, FK}
\end{entidade}

\subsection{Professor}
\begin{entidade}{Professor}
\campo{id\_usuario}{INTEGER}{PK, FK}
\campo{centro}{ENUM}{CCT, CCTA, CBB, CCH, NOT NULL}
\campo{laboratorio}{VARCHAR(100)}{NOT NULL}
\end{entidade}

\subsection{Materia}
\begin{entidade}{Materia}
\campo{id}{SERIAL}{PK}
\campo{nome}{VARCHAR(100)}{NOT NULL}
\campo{codigo\_materia}{VARCHAR(10)}{UNIQUE, NOT NULL}
\end{entidade}

\subsection{Tabela Professor \texorpdfstring{$\times$}{x} Materia (N para N)}
\begin{entidade}{Professor\_x\_Materia}
\campo{id\_professor}{INTEGER}{PK, FK}
\campo{id\_materia}{INTEGER}{PK, FK}
\end{entidade}
\begin{regra}
A chave primária composta \textit{(id\_professor, id\_materia)} impede vincular o mesmo professor duas vezes à mesma matéria.
\end{regra}

\subsection{Aluno}
\begin{entidade}{Aluno}
\campo{id\_usuario}{INTEGER}{PK, FK}
\campo{numero\_matricula}{VARCHAR(20)}{UNIQUE, NOT NULL}
\end{entidade}

\subsection{Bolsista}
\begin{entidade}{Bolsista}
\campo{id\_usuario}{INTEGER}{PK, FK}
\end{entidade}

\subsection{Convite para Aluno}
\begin{entidade}{Convite\_Aluno}
\campo{id}{SERIAL}{PK}
\campo{id\_turma}{INTEGER}{FK, NULL}
\campo{email}{VARCHAR(150)}{NOT NULL}
\campo{convidado\_em}{TIMESTAMP}{NOT NULL}
\campo{status}{ENUM}{pendente, aceitado, expirado, NOT NULL}
\campo{expira\_em}{TIMESTAMP}{NOT NULL}
\campo{convidado\_por}{INTEGER}{FK, NOT NULL}
\campo{numero\_matricula}{VARCHAR(20)}{NULL}
\campo{token\_hash}{VARCHAR(100)}{NOT NULL}
\campo{ultimo\_reenvio\_por}{INTEGER}{FK, NULL}
\campo{exceder\_capacidade}{BOOLEAN}{DEFAULT FALSE, NOT NULL}
\campo{justificativa\_excecao}{TEXT}{NULL}
\campo{aceitado\_por}{INTEGER}{FK, NULL}
\campo{aceitado\_em}{TIMESTAMP}{NULL}
\end{entidade}

\begin{regra}
A combinação \textit{email + id\_turma} não pode possuir mais de um convite \textit{pendente}. Um mesmo e-mail pode possuir convites para turmas diferentes e um aluno pode participar de várias turmas; por isso, o e-mail não é globalmente único nesta entidade.
\end{regra}

\begin{explicacao}
\textit{numero\_matricula} aqui é opcional e \textbf{não é UNIQUE}, pois serve apenas para pré-declarar a matrícula quando o professor a conhece no momento do convite. Ao aceitar o convite, o sistema cria as linhas \textit{Usuario} e \textit{Aluno} correspondentes, copiando \textit{email} e \textit{numero\_matricula} para a nova linha de \textit{Aluno}, onde a restrição \textbf{UNIQUE} é rigorosamente aplicada.
\end{explicacao}

\subsection{Notificação (Unificada)}
\begin{entidade}{Notificacao}
\campo{id}{SERIAL}{PK}
\campo{id\_destinatario}{INTEGER}{FK, NOT NULL}
\campo{papel\_destinatario}{ENUM}{
  \shortstack[l]{%
    Aluno, Professor, Bolsista, Gestor\_Bens\_Patrimoniais,\\
    Gestor\_Almoxarifado, Chefe\_Geral%
  }
}{NOT NULL}
\campo{tipo}{ENUM}{
  \shortstack[l]{%
    ADICIONADO, POST, COMENTARIO, REMOVIDO,\\
    TURMA\_ARQUIVADA, TURMA\_DESARQUIVADA,\\
    REQUISICAO\_BEM, DATA\_DEVOLUCAO\_REAGENTE,\\
    ROTEIRO\_COMPARTILHADO,\\
    REQUISICAO\_EDICAO\_BEM, REQUISICAO\_ADICAO\_BEM,\\
    BEM\_INSERVIVEL,\\
    ENTREGA\_ATRASADA, FRASCOS\_VAZIOS, FRASCOS\_QUEBRADOS,\\
    FRASCOS\_VENCIDOS, FRASCOS\_A\_SEREM\_PESADOS,\\
    FRASCOS\_EM\_QUARENTENA, ESCASSEZ\_ESTOQUE,\\
    AUTO\_ATENDIMENTO\_RETIRADA%
  }
}{NOT NULL}
\campo{id\_quem\_fez\_acao}{INTEGER}{FK, NULL}
\campo{id\_turma}{INTEGER}{FK, NULL}
\campo{quantidade}{INTEGER}{NULL}
\campo{entidade\_alvo}{VARCHAR(40)}{NOT NULL}
\campo{id\_alvo}{VARCHAR(200)}{NOT NULL}
\campo{lida}{BOOLEAN}{DEFAULT FALSE, NOT NULL}
\campo{lida\_em}{TIMESTAMP}{NULL}
\campo{emitida\_em}{TIMESTAMP}{NOT NULL}
\campo{expira\_em}{TIMESTAMP}{ NULL}
\end{entidade}

\begin{explicacao}[Unificação das Notificações]
As quatro tabelas de notificação anteriores (\textit{Notificacao\_para\_Aluno}, \textit{Notificacao\_para\_Professor}, \textit{Notificacao\_para\_Gestor\_de\_bens\_Patrimoniais}, \textit{Notificacao\_para\_Gestor\_de\_Reagentes}) foram consolidadas em uma única entidade \textit{Notificacao}. O campo \textit{papel\_destinatario} indica em qual contexto de papel a notificação deve ser exibida, incluindo Aluno, Professor, Bolsista, Gestor\_Bens\_Patrimoniais, Gestor\_Almoxarifado e Chefe\_Geral.

No Firestore, esta entidade é implementada como sub-coleção \texttt{Usuarios/\{uid\}/Notificacoes}, eliminando a necessidade de queries em múltiplas coleções para usuários multi-role. A ação de ``Limpar tudo'' na interface marca \texttt{lida = true} e \texttt{lida\_em = NOW()} em todas as notificações ativas do usuário, preservando o histórico de auditoria. Os campos \textit{entidade\_alvo} e \textit{id\_alvo} viabilizam a navegação direta (\textit{deep linking}) para o objeto de origem ao clicar no aviso. O campo \textit{expira\_em} é nulo para notificações operacionais persistentes (como \texttt{ESCASSEZ\_ESTOQUE}), que não possuem expiração temporal automática.
\end{explicacao}

\subsection{Tabela Aluno Turma (N para N)}
\begin{entidade}{Aluno\_x\_Turma}
\campo{id\_aluno}{INTEGER}{PK, FK}
\campo{id\_turma}{INTEGER}{PK, FK}
\end{entidade}
\begin{regra}
A chave primária composta \textit{(id\_aluno, id\_turma)} impede matricular o mesmo aluno duas vezes na mesma turma.
\end{regra}

\subsection{Local}
\begin{entidade}{Local}
\campo{id}{SERIAL}{PK}
\campo{predio}{VARCHAR(30)}{NOT NULL}
\campo{andar}{VARCHAR(10)}{NOT NULL}
\campo{sala}{VARCHAR(30)}{NOT NULL}
\campo{criado\_por}{INTEGER}{FK, NOT NULL}
\campo{criado\_em}{TIMESTAMP}{NOT NULL}
\end{entidade}
\begin{regra}
A combinação de prédio, andar e sala deve ser única (\textit{UNIQUE(predio, andar, sala)}).
\end{regra}

\subsection{Post}
\begin{entidade}{Post}
\campo{id}{SERIAL}{PK}
\campo{id\_professor}{INTEGER}{FK, NOT NULL}
\campo{id\_turma}{INTEGER}{FK, NOT NULL}
\campo{id\_roteiro\_experimento}{INTEGER}{FK, NULL}
\campo{titulo}{VARCHAR(150)}{NOT NULL}
\campo{descricao}{TEXT}{NOT NULL}
\campo{criado\_em}{TIMESTAMP}{NOT NULL}
\end{entidade}

\subsection{Comentario}
\begin{entidade}{Comentario}
\campo{id}{SERIAL}{PK}
\campo{id\_post}{INTEGER}{FK, NOT NULL}
\campo{id\_usuario}{INTEGER}{FK, NOT NULL}
\campo{texto}{TEXT}{NOT NULL}
\campo{criado\_em}{TIMESTAMP}{NOT NULL}
\campo{moderado}{BOOLEAN}{DEFAULT FALSE, NOT NULL}
\campo{motivo\_moderacao}{TEXT}{NULL}
\campo{moderado\_por}{INTEGER}{FK, NULL}
\end{entidade}

\subsection{Registro de Auditoria}
\begin{entidade}{Registro\_de\_Auditoria}
\campo{id}{SERIAL}{PK}
\campo{id\_usuario}{INTEGER}{FK, NOT NULL}
\campo{acao}{VARCHAR(200)}{NOT NULL}
\campo{tipo\_entidade\_sofre\_acao}{ENUM}{
  \shortstack[l]{%
    TURMA, ALUNO, PROFESSOR, ALMOXARIFADO, RESUMO\_REAGENTE,\\
    ESPECIFICACAO\_REAGENTE, LOTE, FRASCO\_REAGENTE,\\
    EMPRESTIMO\_REAGENTE, BEM\_PATRIMONIAL, LOCAL, MATERIA,\\
    COMENTARIO, POST, ROTEIRO, REQUISICAO\_ADICAO\_BEM,\\
    REQUISICAO\_EDICAO\_BEM, USUARIO%
  }
}{NOT NULL}
\campo{id\_do\_objeto\_da\_entidade}{VARCHAR(200)}{NOT NULL}
\campo{acao\_feita\_em}{TIMESTAMP}{NOT NULL}
\campo{metadata}{JSONB}{NOT NULL}
\end{entidade}

\begin{explicacao}
\textit{acao} e \textit{id\_do\_objeto\_da\_entidade} tinham o tamanho \texttt{200} escrito como se fosse um selo de restrição em vez de parte do tipo; corrigido para \texttt{VARCHAR(200)}. O valor \texttt{REAGENTE} do enum foi substituído por valores granulares.
\textbf{Nota sobre Impressão:} A reimpressão avulsa (2ª via) de um frasco já cadastrado é auditada diretamente aqui (\texttt{tipo = FRASCO\_REAGENTE}, \texttt{acao = REIMPRESSAO\_ETIQUETA}) em vez de usar a entidade \textit{Impressao\_Etiqueta\_Frasco}.
\end{explicacao}
\newpage

\subsection{Especificação de Integridade do Modelo 3FN}
Esta subseção fecha as propriedades lógicas que devem permanecer verdadeiras antes de qualquer denormalização específica do Firestore.

\subsubsection{Cardinalidades}
\begin{longtable}{>{\raggedright\arraybackslash}p{0.36\textwidth} >{\centering\arraybackslash}p{0.14\textwidth} >{\raggedright\arraybackslash}p{0.42\textwidth}}
\toprule
\textbf{Relacionamento} & \textbf{Card.} & \textbf{Regra}\\
\midrule
\endfirsthead
\toprule
\textbf{Relacionamento} & \textbf{Card.} & \textbf{Regra} (continuação)\\
\midrule
\endhead
\bottomrule
\endfoot
Usuario -- cada papel específico & 0..1 por papel & Um usuário pode ter no máximo uma linha em cada tabela de papel; várias tabelas de papel podem estar associadas ao mesmo usuário conforme as combinações permitidas.\\
Gestor Almoxarifado -- Almoxarifado & N:N & A mesma associação não pode ocorrer duas vezes.\\
Resumo Bem -- Bem Patrimonial & 1:N & Cada bem pertence a um resumo.\\
Local -- Bem Patrimonial & 1:N & Cada bem possui um local cadastrado.\\
Resumo Reagente -- Especificação & 1:N & Um resumo agrupa especificações; cada especificação pertence a um resumo.\\
Especificação -- Substância & N:N & A composição liga múltiplas substâncias a múltiplas especificações.\\
Especificação -- Lote & 1:N & Um lote pertence a uma única especificação.\\
Lote -- Frasco & 1:N opcional & Um frasco pode ter lote NULL quando o lote não for conhecido.\\
Almoxarifado -- Frasco & 1:N & Cada frasco está em exatamente um almoxarifado.\\
Frasco -- Empréstimo & 1:N & Histórico completo; no máximo um empréstimo ativo por frasco.\\
Turma -- Aluno & N:N & Alunos podem participar de várias turmas (via sub-coleção \texttt{Turma/\{id\}/Alunos}).\\
Turma -- Post & 1:N & Cada post pertence a uma turma e ao professor autor.\\
Post -- Comentário & 1:N & Cada comentário pertence a um post.\\
Professor -- Matéria & N:N & A associação não pode ser duplicada.\\
Roteiro -- Professor & N:N & O mesmo roteiro não pode ser compartilhado duas vezes com o mesmo professor.\\
\bottomrule
\end{longtable}

\subsubsection{Nullabilidade}
Todo atributo é \textbf{NOT NULL} por padrão. Um campo só pode ser NULL quando isso representa uma ausência legítima do domínio. As exceções são:
\begin{itemize}
    \item foto de perfil;
    \item documento de baixa antes da baixa patrimonial;
    \item campos de resposta de requisição antes da resposta;
    \item campos opcionais de especificação (fabricante, código do fabricante, FDS/FISPQ e demais dados não disponíveis);
    \item CAS e fórmula química quando não cadastrados;
    \item frequência de pesagem quando o reagente não requer pesagem frequente;
    \item lote e dados de validade/pesagem quando desconhecidos ou ainda não aplicáveis;
    \item validade\_efetiva antes da abertura quando sua definição depende do prazo após abertura;
    \item dados de devolução de empréstimo antes da devolução;
    \item id\_local da linha agregada ``Geral'' no resumo patrimonial diário;
    \item id\_turma da notificação de professor quando o tipo não é COMENTARIO;
    \item atributos de composição cuja ausência é necessária para representar faixa ou notação textual do fabricante.
\end{itemize}

\subsubsection{Chaves e constraints}
\begin{itemize}
    \item Usuario.email, Bem\_Patrimonial.numero\_patrimonio, Frasco\_Reagente.codigo\_frasco, Turma.codigo\_turma, Materia.codigo\_materia e CAS não nulo são únicos.
    \item As tabelas associativas usam chave primária composta por suas duas FKs.
    \item Composicao\_Reagente usa UNIQUE(id\_especificacao\_reagente, id\_substancia\_quimica).
    \item Lote usa UNIQUE(id\_especificacao\_reagente, nome\_fornecedor, numero\_lote), evitando a identificação de dois lotes equivalentes do mesmo fornecedor e produto.
    \item Convite\_Aluno permite o mesmo e-mail em turmas diferentes, mas impede mais de um convite PENDENTE para a mesma combinação de e-mail normalizado e turma; para convite global (\textit{id\_turma} NULL), impede mais de um convite PENDENTE para o mesmo e-mail normalizado. No PostgreSQL, isso é representado por dois índices únicos parciais: \texttt{UNIQUE(email\_normalizado, id\_turma) WHERE id\_turma IS NOT NULL AND status = 'pendente'} e \texttt{UNIQUE(email\_normalizado) WHERE id\_turma IS NULL AND status = 'pendente'}. No Firestore, a Cloud Function deve realizar a mesma validação dentro de transação.
    \item Requisicao\_Edicao\_Bem\_Patrimonial permite no máximo uma requisição PENDENTE por bem.
    \item Todos os valores quantitativos de massa, volume, capacidade, contagens, mês, dia de frequência e quantidades agregadas devem obedecer aos limites não negativos ou positivos definidos pelas regras do domínio.
    \item Datas que representam intervalos devem respeitar sua ordem temporal; exemplos: fabricação <= validade e devolução >= retirada.
\end{itemize}

\subsubsection{Regras químicas fechadas}
\begin{itemize}
    \item \textbf{PURA} e \textbf{MISTURA} são tipos de substância do resumo. MISTURA exige composição na especificação; PURA não depende de múltiplas linhas de composição.
    \item A concentração, quando aplicável, é atributo da composição da especificação, e não do resumo.
    \item A densidade é atributo da especificação. É obrigatória e positiva para LIQUIDO e pode ser NULL para SOLIDO.
    \item LIQUIDO usa ml como unidade operacional; SOLIDO usa g. O par estado/unidade deve ser sempre coerente.
    \item Duas especificações com o mesmo nome podem coexistir sob o mesmo resumo quando diferirem em concentração, densidade, fabricante ou código do produto e representarem o mesmo item agrupador. Diferença de especificação não implica automaticamente novo resumo.
    \item Gases estão fora do escopo da V1. Não se deve misturar sua futura regra de pressão com os campos de pesagem de sólidos e líquidos.
\end{itemize}

\subsubsection{Máquina de estados e invariantes do frasco}
\begin{longtable}{>{\raggedright\arraybackslash}p{0.34\textwidth} >{\raggedright\arraybackslash}p{0.58\textwidth}}
\toprule
\textbf{Estado/ação} & \textbf{Invariante}\\
\midrule
\endfirsthead
\toprule
\textbf{Estado/ação} & \textbf{Invariante} (continuação)\\
\midrule
\endhead
\bottomrule
\endfoot
FECHADO + DISPONIVEL & Pode ser retirado ou aberto.\\
ABERTO + DISPONIVEL & Pode ser retirado ou ajustado por pesagem.\\
EMPRESTADO & Deve possuir exatamente um empréstimo ativo EM\_USO ou ATRASADO.\\
VAZIO/QUEBRADO & Não pode estar EMPRESTADO e deve aparecer como pendente de descarte.\\
QUARENTENA & Exige detalhe de motivo e bloqueia nova retirada.\\
VENCIDO + uso\_vencido\_autorizado = false & Exibe PENDENTE DE DESCARTE quando não estiver em quarentena; bloqueia nova retirada.\\
VENCIDO + uso\_vencido\_autorizado = true & Pode permanecer DISPONIVEL para uso excepcional após confirmação explícita no empréstimo.\\
DESCARTADO & É terminal; não retorna a DISPONIVEL.\\
Devolução & Só pode ocorrer para empréstimo EM\_USO ou ATRASADO e aplica tolerância híbrida Q06: ganho tolerado produz consumo zero e evento AJUSTE; ganho acima da tolerância bloqueia.\\
\bottomrule
\end{longtable}

\subsubsection{Fechamento dos estados automáticos}
\textit{vencido} é definido pelo servidor a partir da validade efetiva. Na abertura do frasco, \textit{validade\_efetiva} é recalculada conforme a regra de validade pós-abertura. \textit{ATRASADO} é promovido pelo job servidor quando a data de devolução prevista é ultrapassada. Nenhum desses estados pode ser inventado pelo frontend. O vencimento também é recalculado sincronicamente em operações críticas de cadastro e devolução, para que o resultado não dependa de o job diário já ter executado.

\subsection{Checklist de Fechamento do Domínio e do Modelo}
Os onze pontos que precisavam ser resolvidos antes da modelagem física foram consolidados da seguinte forma:
\begin{enumerate}
    \item \textbf{Matriz completa de permissões}: formalizada na seção de Stakeholders e referenciada como fonte normativa de autorização.
    \item \textbf{Chefe Geral e multi-role}: Chefe Geral é exclusivo por conta; eventual dupla função institucional exige contas independentes.
    \item \textbf{Obrigatoriedade de campos}: formalizada por contexto, estado, tipo de substância e operação.
    \item \textbf{Densidade}: reposicionada semanticamente como propriedade de Especificacao\_Reagente e usada nos cálculos sem duplicação no resumo.
    \item \textbf{Máquina de estados}: estados do frasco e transições de retirada, devolução, quarentena, vencimento, esvaziamento, quebra e descarte foram explicitados.
    \item \textbf{NULL/NOT NULL}: todo campo passa a ter semântica explícita, com NULL restrito a ausências legítimas do domínio.
    \item \textbf{Cardinalidades}: relacionamentos 1:N, N:N e especializações de papel foram consolidados.
    \item \textbf{Constraints}: unicidades, chaves compostas, regras temporais e limites quantitativos foram formalizados, indicando quais precisarão ser simulados por transação no Firestore.
    \item \textbf{Métricas e unidades}: limiar de escassez passa a significar quantidade de frascos; massa de balança e consumo convertido permanecem conceitos distintos.
    \item \textbf{Gases}: explicitamente fora do escopo da V1, com a extensão futura isolada para não contaminar o modelo atual.
    \item \textbf{Vencimento e destino}: cadastro, abertura e devolução distinguem explicitamente quarentena, pendência de descarte e uso vencido autorizado.
\end{enumerate}

\begin{explicacao}[Critério de passagem para a modelagem Firestore]
O modelo 3FN é considerado fechado quando cada requisito funcional consegue ser rastreado até entidade, atributo, constraint e fluxo correspondente sem depender de uma suposição não documentada. A seção seguinte pode então introduzir apenas as denormalizações necessárias às consultas do Firestore, sem alterar a semântica do modelo relacional.
\end{explicacao}

\begin{regra}[DP-C02: moderação]
Autor vê o original marcado como moderado; demais alunos e bolsistas veem aviso institucional. Professor responsável e Chefe Geral veem o original e histórico necessário à moderação/auditoria. Moderação exige motivo e moderador; proteção ocorre no servidor e nas Rules (seção 11).
\end{regra}
\begin{regra}[DP-C03: segredo de convite]
Gerar token com CSPRNG de 32 bytes e armazenar somente SHA-256 do token aleatório. Comparar em tempo constante no backend, validar expiração de sete dias, identidade verificada e consumo único transacional. Reenvio substitui hash e expiração no mesmo documento e registra o último operador; link anterior deixa de funcionar. No Firestore, docId determinístico é derivado no servidor por HMAC-SHA-256 com segredo do servidor sobre serialização inequívoca de contexto turma/global e e-mail normalizado. Não expor hash simples de e-mail previsível. HMAC é pseudonimização; documento e e-mail continuam restritos. Nunca registrar token em logs, auditoria ou banco.
\end{regra}

\begin{regra}[Q06: massa, consumo e tolerância]
Massa consumida = $\max(0, peso_{saida}-peso_{retorno})$ em g. Medida utilizada é não negativa (g para sólido, mL via densidade para líquido). Medida usada no frasco acumula apenas gramas. Tolerância normal: $\max(1\,g,0{,}005\,peso_{saida})$; higroscópico: $\max(2\,g,0{,}02\,peso_{saida})$. A fórmula de tolerância higroscópica e de variação instrumental incide estritamente sobre o \textbf{Peso Bruto de Saída ($peso_{saida}$)}. Ganho dentro da tolerância preserva peso físico real, consumo zero e evento AJUSTE/ganho\_massa\_higroscopia. Retorno abaixo da tara ($peso_{retorno} < peso_{vazio}$): diferença até 5 g exige confirmação de esgotamento para VAZIO e ajuste auditado; diferença maior com produto restante bloqueia e exige recalibração formal de tara com justificativa. Não converter ajuste em consumo negativo.
\end{regra}
\begin{regra}[Q14 e Q04: rastreabilidade na retirada]
Autoatendimento é calculado pelo servidor: só o único gestor ativo do almoxarifado pode retirar para si, com justificativa obrigatória e notificação auditável à chefia; se houver outro gestor ativo, bloquear. Para vencido ou validade desconhecida liberado com termo nas finalidades PESQUISA\_TCC\_POS ou ESTUDO\_DEGRADACAO\_RESIDUOS, exigir justificativa metodológica de 20--2000 caracteres, versão do TCR, instante servidor, UID autenticado do retirante que aceitou e referência ao Registro\_de\_Auditoria correspondente. O backend valida aceite vinculado ao frasco, finalidade e operação, nunca declaração do gestor em nome de terceiro. Sem essas condições, rejeitar. TCR registra ciência e responsabilidade, sem substituir regras institucionais de segurança química. Quarentena e pendência de descarte continuam bloqueadas.
\end{regra}
\begin{regra}[Reclassificação patrimonial]
Novo nome individual exige resolver/criar resumo alvo, preencher novo\_id\_resumo\_bem\_patrimonial e alterar só o vínculo do bem solicitado; não renomear resumo compartilhado. Versao\_bem\_origem captura a versão observada e impede aprovação sobre edição concorrente.
\end{regra}
